\section{Conclusion}\label{sec:conclusion}

Calibrating a variant effect predictor variant by variant lets well-studied families set its error rate: on ProteinGym, Mondrian split conformal covers only \sOneBtwoPath{} of pathogenic variants per family for \sOneBtwoNBelow{} predictors, without any shift. Treating families as the exchangeable unit removes this failure, and conditioning a family-weighted quantile regression on a cross-fitted density ratio guarantees coverage under family exchangeability and when the estimated ratio is correct. The empirical gain is in coverage and appears mainly under descriptor shift (S2); on orphan families and held-out embedding regions ProtCOVER matches group weighting. On proteins outside the benchmark, a pre-registered test with two predictors found variant-level calibration failing again and every method short of the target where the predictors were less accurate and the longest proteins lacked calibration support, neither of which the label-free ratio flags; conditioning narrowed the shortfall, within the margin of group weighting, without closing it. Descriptors that track a predictor's accuracy without labels, and calibration support for proteins unlike any calibration protein, remain open.
